% GENERATED FILE. Edit canonical lessons, then run npm run generate:books.
% canonical-source-hash: fnv1a64:b451d0cda8ee4195
% canonical-lessons: MR-C325-raja, MR-C325-rani, MR-C325-rajvada, MR-C325-darbar, MR-C325-rajya-karne

\chapter{A King and His Court}
\label{ch:mr-a-king-and-his-court}
\hlchaptermodality{325}

\begin{chapteropening}
\textbf{By the end of this chapter:} \emph{I can tell about a king, a queen and their palace.}

Complete the last lesson of chapter 325: I can tell about a king, a queen and their palace.
\end{chapteropening}

\section[\texorpdfstring{\mr{राजा} (rājā)}{राजा (rājā)}]{\mr{राजा} (rājā) — a king}

\label{lesson:MR-C325-raja}



\begin{quote}
Before the new one: say the Marathi for to go off, then the Marathi for to compose.
\end{quote}

\subsection*{You'll want to know: \mr{राजा}}
\textbf{\mr{राजा}} — \emph{rājā} — \textquotedblleft{}a king\textquotedblright{}.

\begin{quote}
\mr{एकेकाळी} \mr{या} \mr{प्रदेशात} \mr{एक} \mr{दयाळू} \mr{राजा} \mr{होता}.
\end{quote}

\begin{quote}
\emph{ekekāḷī yā pradeśāt ek dayāḷū rājā hotā.}
\end{quote}

\begin{quote}
Once upon a time there was a kind king in this region.
\end{quote}

Masculine. Plural: \mr{राजे}, \emph{rāje}.

Plural \mr{राजे}.

\begin{cousinweb}
\textbf{\mr{राजा}} (\emph{rājā}) is a tatsama, taken from Sanskrit \emph{rājan}, 'a king', in its nominative form \emph{rājā} (Monier-Williams).
\end{cousinweb}

\begin{grammarlens}[title={What you've built}]
One more word for telling what happened.
\end{grammarlens}

\subsection*{Your turn}
\begin{itemize}
\raggedright
  \item \emph{Say it:} \emph{rājā}
  \item \emph{Say it:} \emph{rājā}, once more
  \item \emph{Say it:} \emph{ekekāḷī yā pradeśāt ek dayāḷū rājā hotā.} — Once upon a time there was a kind king in this region.
  \item \emph{Recall:} say the Marathi for to go off, then the Marathi for to compose, then say \emph{rājā} again
\end{itemize}

\begin{tcolorbox}[breakable,colback=teal!4,colframe=teal!35!black,title={Before you move on}]
What does \mr{राजा} mean? (\textquotedblleft{}A king\textquotedblright{}.) Say it once more.
\end{tcolorbox}
\section[\texorpdfstring{\mr{राणी} (rāṇī)}{राणी (rāṇī)}]{\mr{राणी} (rāṇī) — a queen}

\label{lesson:MR-C325-rani}



\begin{quote}
Before the new one: say the Marathi for to compose, then the Marathi for a king.
\end{quote}

\subsection*{You'll want to know: \mr{राणी}}
\textbf{\mr{राणी}} — \emph{rāṇī} — \textquotedblleft{}a queen\textquotedblright{}.

\begin{quote}
\mr{राणीने} \mr{गरीब} \mr{लोकांसाठी} \mr{विहिरी} \mr{बांधल्या}.
\end{quote}

\begin{quote}
\emph{rāṇīne garīb lokāṁsāṭhī vihirī bāndhalyā.}
\end{quote}

\begin{quote}
The queen built wells for poor people.
\end{quote}

Feminine.

Plural \mr{राण्या}.

\begin{cousinweb}
\textbf{\mr{राणी}} (\emph{rāṇī}) is a tadbhava: CDIAL lists Marathi \emph{rāṇī}, 'a queen', under Sanskrit \emph{rājñī}, 'a queen' (CDIAL 10692).
\end{cousinweb}

\begin{grammarlens}[title={What you've built}]
One more word for telling what happened.
\end{grammarlens}

\subsection*{Your turn}
\begin{itemize}
\raggedright
  \item \emph{Say it:} \emph{rāṇī}
  \item \emph{Say it:} \emph{rāṇī}, once more
  \item \emph{Say it:} \emph{rāṇīne garīb lokāṁsāṭhī vihirī bāndhalyā.} — The queen built wells for poor people.
  \item \emph{Recall:} say the Marathi for to compose, then the Marathi for a king, then say \emph{rāṇī} again
\end{itemize}

\begin{tcolorbox}[breakable,colback=teal!4,colframe=teal!35!black,title={Before you move on}]
What does \mr{राणी} mean? (\textquotedblleft{}A queen\textquotedblright{}.) Say it once more.
\end{tcolorbox}
\section[\texorpdfstring{\mr{राजवाडा} (rājvāḍā)}{राजवाडा (rājvāḍā)}]{\mr{राजवाडा} (rājvāḍā) — a palace}

\label{lesson:MR-C325-rajvada}



\begin{quote}
Before the new one: say the Marathi for a king, then the Marathi for a queen.
\end{quote}

\subsection*{You'll want to know: \mr{राजवाडा}}
\textbf{\mr{राजवाडा}} — \emph{rājvāḍā} — \textquotedblleft{}a palace\textquotedblright{}.

\begin{quote}
\mr{जुन्या} \mr{राजवाड्यात} \mr{आता} \mr{संग्रहालय} \mr{आहे}.
\end{quote}

\begin{quote}
\emph{junyā rājvāḍyāt ātā saṅgrahālay āhe.}
\end{quote}

\begin{quote}
There is now a museum in the old palace.
\end{quote}

Masculine.

The oblique is \mr{राजवाड्या}-: \mr{राजवाड्यात}, in the palace.

\begin{cousinweb}
\textbf{\mr{राजवाडा}} (\emph{rājvāḍā}) is built from \emph{rāj-}, 'royal', and \emph{vāḍā}, 'a mansion, an enclosure', which goes back to Sanskrit \emph{vāṭa}, 'an enclosure' (CDIAL 11480).
\end{cousinweb}

\begin{grammarlens}[title={What you've built}]
One more word for telling what happened.
\end{grammarlens}

\subsection*{Your turn}
\begin{itemize}
\raggedright
  \item \emph{Say it:} \emph{rājvāḍā}
  \item \emph{Say it:} \emph{rājvāḍā}, once more
  \item \emph{Say it:} \emph{junyā rājvāḍyāt ātā saṅgrahālay āhe.} — There is now a museum in the old palace.
  \item \emph{Recall:} say the Marathi for a king, then the Marathi for a queen, then say \emph{rājvāḍā} again
\end{itemize}

\begin{tcolorbox}[breakable,colback=teal!4,colframe=teal!35!black,title={Before you move on}]
What does \mr{राजवाडा} mean? (\textquotedblleft{}A palace\textquotedblright{}.) Say it once more.
\end{tcolorbox}
\section[\texorpdfstring{\mr{दरबार} (darbār)}{दरबार (darbār)}]{\mr{दरबार} (darbār) — a royal court}

\label{lesson:MR-C325-darbar}



\begin{quote}
Before the new one: say the Marathi for a queen, then the Marathi for a palace.
\end{quote}

\subsection*{You'll want to know: \mr{दरबार}}
\textbf{\mr{दरबार}} — \emph{darbār} — \textquotedblleft{}a royal court\textquotedblright{}.

\begin{quote}
\mr{राजाच्या} \mr{दरबारात} \mr{लोक} \mr{तक्रारी} \mr{घेऊन} \mr{यायचे}.
\end{quote}

\begin{quote}
\emph{rājācyā darbārāt lok takrārī gheūn yāyce.}
\end{quote}

\begin{quote}
People used to come to the king's court with complaints.
\end{quote}

Masculine.

\mr{दरबारात}, at court.

\begin{cousinweb}
\textbf{\mr{दरबार}} (\emph{darbār}) is a loan from Persian \emph{darbār}, 'a royal court, an audience hall'.
\end{cousinweb}

\begin{grammarlens}[title={What you've built}]
One more word for telling what happened.
\end{grammarlens}

\subsection*{Your turn}
\begin{itemize}
\raggedright
  \item \emph{Say it:} \emph{darbār}
  \item \emph{Say it:} \emph{darbār}, once more
  \item \emph{Say it:} \emph{rājācyā darbārāt lok takrārī gheūn yāyce.} — People used to come to the king's court with complaints.
  \item \emph{Recall:} say the Marathi for a queen, then the Marathi for a palace, then say \emph{darbār} again
\end{itemize}

\begin{tcolorbox}[breakable,colback=teal!4,colframe=teal!35!black,title={Before you move on}]
What does \mr{दरबार} mean? (\textquotedblleft{}A royal court\textquotedblright{}.) Say it once more.
\end{tcolorbox}
\section[\texorpdfstring{\mr{राज्य} \mr{करणे} (rājya karṇe)}{राज्य करणे (rājya karṇe)}]{\mr{राज्य} \mr{करणे} (rājya karṇe) — to reign}

\label{lesson:MR-C325-rajya-karne}



\begin{quote}
Before the new one: say the Marathi for a palace, then the Marathi for a royal court.
\end{quote}

\subsection*{You'll want to know: \mr{राज्य} \mr{करणे}}
\textbf{\mr{राज्य} \mr{करणे}} — \emph{rājya karṇe} — \textquotedblleft{}to reign\textquotedblright{}.

\begin{quote}
\mr{त्या} \mr{राजाने} \mr{चाळीस} \mr{वर्षं} \mr{राज्य} \mr{केलं}.
\end{quote}

\begin{quote}
\emph{tyā rājāne cāḷīs varṣaṁ rājya kelaṁ.}
\end{quote}

\begin{quote}
That king reigned for forty years.
\end{quote}

Past, with -\mr{ने} (he …): \mr{त्याने} \mr{राज्य} \mr{केलं}, \emph{tyāne rājya kelaṁ}.

The land ruled takes -\mr{वर}: \mr{देशावर} \mr{राज्य} \mr{करणे}.

\begin{cousinweb}
\textbf{\mr{राज्य} \mr{करणे}} (\emph{rājya karṇe}) joins \emph{rājya}, a tatsama taken from Sanskrit \emph{rājya}, 'kingship, a kingdom' (Monier-Williams), to \emph{karṇe}, 'to do'.
\end{cousinweb}

\begin{grammarlens}[title={What you've built}]
One more word for telling what happened.
\end{grammarlens}

\subsection*{Your turn}
\begin{itemize}
\raggedright
  \item \emph{Say it:} \emph{rājya karṇe}
  \item \emph{Say it:} \emph{rājya karṇe}, once more
  \item \emph{Say it:} \emph{tyā rājāne cāḷīs varṣaṁ rājya kelaṁ.} — That king reigned for forty years.
  \item \emph{Recall:} say the Marathi for a palace, then the Marathi for a royal court, then say \emph{rājya karṇe} again
\end{itemize}

\begin{tcolorbox}[breakable,colback=teal!4,colframe=teal!35!black,title={Before you move on}]
What does \mr{राज्य} \mr{करणे} mean? (\textquotedblleft{}To reign\textquotedblright{}.) Say it once more.
\end{tcolorbox}
